\section{Projects}
    \vspace{-5pt}
    \resumeSubHeadingListStart
      \resumeProjectHeading
          {\textbf{LambdaMailer} $|$ \emph{Node.js, Amazon Lambda, Amazon SES}  \hspace{0.2cm}
          \href{https://github.com/RudradevArya/Lambda-Mailer-Backend}{\Large\faGithub}}{February 2024}
          \resumeItemListStart
            \resumeItem{Leveraged \textbf{AWS Beanstalk} for scalable backend deployment and \textbf{Vercel} for frontend with which users can upload a CSV with relevant target mail IDs and other information}
            \resumeItem{Acheived a success rate of \textbf{99.2\%} for \textbf{50,000} emails at just \textbf{\$40/month} at maximum throughput of \textbf{200 mails/sec} and at maximum capacity}
            \resumeItem{Utilized \textbf{React} and \textbf{React Router} for a user friendly interface and \textbf{OAuth 2.0} protocol for secure authentication and authorization and minimal}
            \resumeItem{Designed the project as a microservice, communicating via APIs built with \textbf{Axios} and \textbf{Express.js}}
          \resumeItemListEnd
          \vspace{-13pt}
          \resumeProjectHeading
          {\textbf{EventU} $|$ \emph{Next.js, Appwrite, TypeScript} \hspace{0.2cm}  
          \href{https://github.com/RudradevArya/EventU}{\Large\faGithub}} {August 2023}         
          \resumeItemListStart
            \resumeItem{Engineered a complex event management system leveraging \textbf{NoSQL} databases; streamlined unstructured data processing and achieved a \textbf{40\%} reduction in data retrieval time, enhancing overall event coordination efficiency}
            \resumeItem{Designed the user interface using HTML, \textbf{TailwindCSS}, and \textbf{Next.js}, enhancing the platform’s usability.}
            \resumeItem{Leveraged \textbf{Appwrite} for authentication, multiple sign-in methods, and efficient storage and file management.}
            \resumeItem{Utilized Next.js, an open-source React-based framework, to achieve server-side rendering and enhanced \textbf{SEO}. Deployed the application on \textbf{Vercel’s} frontend cloud infrastructure for optimal performance.}
          \resumeItemListEnd 
          \vspace{-13pt}
          \resumeProjectHeading
          {\textbf{Elevator Waiting Time Estimation and Simulation} $|$ \emph{Arduino PoC, Python, Java, C++}  \hspace{0.2cm}
          \href{https://github.com/RudradevArya/Elevator-Simulation-nd-Estimation}{\large\faGithub} and 
          \href{https://github.com/RudradevArya/Elevator-Waiting-Time-Estimation-System}{\large\faGithub}}{Jan 2023}
          \resumeItemListStart
            \resumeItem{Implemented Various \textbf{Algorithms} using \textbf{Java} like Round Robin, Up Peak, Zoning etc to effectively simulate \textbf{Scheduling Algorithms} in an elevator}
            \resumeItem{Created a \textbf{Proof Of Concept(PoC)} model using Arduino and C++ to demonstrate real world effectiveness of Scheduling algorithms and to effectively \textbf{estimate} waiting time for an elevator }
            
          \resumeItemListEnd
          \vspace{-13pt}
      
          \resumeProjectHeading
          {\textbf{Minecraft Replica using Python } $|$ \emph{Python, Pyglet, Linear Algebra} \hspace{0.2cm}
          \href{https://github.com/RudradevArya/Python-Minecraft}{\Large\faGithub}}{April 2023}
          \resumeItemListStart
            \resumeItem{Utilized \textbf{Pyglet} a Python multimedia library to create a \textbf{128x128} Block replica of the original Minecraft game}
            \resumeItem{Leverage basic \textbf{Linear Algebra} to understand various game development techniques such as cursor and player movement correlation}
            \resumeItem{Implemented \textbf{object-oriented programming} practices such as abstraction to abstract away user controls and player movements}
            \resumeItem{Implemented \textbf{CUDA} to parallelize the workload of the game improving the FPS from \textbf{50 FPS} on CPU to over \textbf{270 FPS} on a mid tier GPU}
          \resumeItemListEnd 
    \resumeSubHeadingListEnd
\vspace{-15pt}
